\section{Eliminating the activation-order search}
\label{sec:factorial-free}

The ordered algorithm makes causality explicit, but enumerating all $k!$
orders is unnecessary. A single distinguished block can encode the
prerequisites for which no initial-seed witness is requested. Forward
chaining then determines whether the chosen outputs can justify the other
activations. This yields the strongest bound proved in this article.

\begin{theorem}[Output choices and one distinguished partition block]
\label{thm:factorial-free}
The triple-unit forest problem is decidable, with a successful tree and
peeling certificate when one exists, in
\[
                 3^k B_{2k+1}\operatorname{poly}(N)
                 =2^{O(k\log(k+2))}\operatorname{poly}(N)
\]
bit operations. For arbitrary active exceptional supports $H_i$, with
unit-owner sets $U_i$ and $Q=\sum_i(|H_i|-1)$, the bound is
\[
                 \left(\prod_i |U_i|\right)B_{Q+1}
                 \operatorname{poly}(N).
\]
Both statements hold with an arbitrary polynomial-oracle matroid in place
of the forest independence constraint.
\end{theorem}

\subsection{The reduced search space}

Fix one unit output $o_i\in H_i$ for each exceptional column. No event
order is selected. Retain only its implication
$H_i\setminus\{o_i\}\to o_i$ in the nonlinear part of this branch.
Write $\mathcal C_o$ for closure under these chosen implications and the
ordinary digraph $P$.

As before, let $Q_1,\ldots,Q_p$ be the source SCCs of $P$ containing
neither a free seed nor any chosen output. Set
\[
 A_j=Q_j\cap E,\qquad U=Z\cup\bigcup_j Q_j,
 \qquad
 X=\left(\bigcup_i(H_i\setminus\{o_i\})\right)
                         \setminus\Reach_P(U).
\]
Repeated occurrences of the same prerequisite row now need only one
index. Let $q=|X|\le Q$, and for $x\in X$ put
\[
                         D_x=\{e\in E:x\in\Reach_P(\{e\})\}.
\]
Every $D_x$ is disjoint from every $A_j$, by exactly the separation
argument used in Lemma~\ref{lem:fpt-demand-separation}.

Introduce a new symbol $\star\notin R$ and enumerate set partitions
$\pi$ of $X\cup\{\star\}$. Denote by $B_\star$ the unique block
containing $\star$, and define
\[
                         S_\pi=X\setminus B_\star.
\]
The rows of $S_\pi$ are required to be reached from initial forest
witnesses. Rows in $B_\star\setminus\{\star\}$ are \emph{undemanded}:
they are allowed to activate in any legal way, even initially. The
construction never requires them to remain inactive.

For this partition, first perform the purely logical test
\begin{equation}
                       \mathcal C_o(U\cup S_\pi)=R.
                       \label{eq:output-closure-test}
\end{equation}
If it fails, discard the partition. If it succeeds, seek independent
distinct representatives for the groups
\begin{equation}
 A_1,\ldots,A_p,\qquad
 K_B=\bigcap_{x\in B}D_x
       \quad(B\in\pi,\ B\ne B_\star).
 \label{eq:output-groups}
\end{equation}
The same cloned graphic/partition intersection solves this subproblem.
Empty intersections are immediate obstructions.

\begin{proof}[Proof of Theorem~\ref{thm:factorial-free}]
For soundness, an independent representative set $F$ activates every
mandatory source SCC, and therefore all of $U$. Its representative for
$K_B$ reaches every row of $B$, so all of $S_\pi$ is active. The closure
test then supplies a legal activation sequence for every row. Extending
$F$ to a spanning tree preserves success. A literal matrix peeling replay
produces the positive certificate.

For completeness, take a successful independent seed $F$ and an actual
complete activation schedule. For every exceptional column used to
activate a new row, choose that row as its output $o_i$. Choose any unit
output for every unused column. Each chosen mandatory source SCC contains
an element of $F$, since it has no incoming ordinary arc, free seed, or
chosen nonlinear output.

Let $D=\Reach_P(Z\cup F)$ and $S=X\cap D$. We claim that
$\mathcal C_o(U\cup S)=R$. Replay the useful exceptional events of
the actual schedule. Each prerequisite of such an event either belongs
to $D$ or is ordinarily reachable from an earlier useful output. In the
first case it lies in $\Reach_P(U)$ or in $S$, by the definition of $X$;
in the second it is available by induction. Every useful chosen output
is consequently justified in the restricted system.

After all useful events, every prerequisite of an unused exceptional
column is available as well. Indeed, completion of the original schedule
places that prerequisite in $D$ or in ordinary reach from a useful
output, and the same argument applies. Thus the arbitrary chosen outputs
of unused columns can now be added. Every source SCC of $P$ contains a
free seed, a mandatory seed, or one of these chosen outputs. All rows
are therefore reached, proving the claim. This argument includes chosen
outputs that were already in the initial seed; it does not presume that
every selected output was a useful event in the original schedule.

For each $x\in S$, choose $f_x\in F\cap D_x$. Such a witness exists
because $x\in D$ and $x\notin\Reach_P(U)$, with $Z\subseteq U$.
Group the elements of $S$ by equality of their chosen witnesses, and put
$\{\star\}\cup(X\setminus S)$ in the distinguished block. Choose a
representative from $F\cap A_j$ for each mandatory group. The demand
separation ensures that all these group representatives are distinct;
they form an independent subset of $F$. The constructed partition passes
both the closure test and the independent-transversal test.

There are $\prod_i|U_i|$ output assignments. For each, at most $Q$
distinct prerequisites remain, so there are at most $B_{Q+1}$ partitions.
Everything inside one branch is polynomial in the input length. This
proves both bounds and their matroid extension.
\end{proof}

Equivalently, first choose $S\subseteq X$, then partition its demands.
The exact count
\[
                    \sum_{s=0}^{q}\binom qs B_s=B_{q+1}
\]
explains why the distinguished block absorbs both choices. In particular,
there is no additional factor $2^q$ multiplying $B_q$. The two exact
strategies give the combined theoretical bound
\[
  3^k\min\{k!B_{2k},B_{2k+1}\}\operatorname{poly}(N).
\]
This improves the explicit parameter factor but leaves its
$2^{O(k\log k)}$ asymptotic class unchanged. It is not a single-exponential
algorithm in $k$.

\subsection{Rank, weights, and independent negative replay}

Under the full raw parity hypothesis, a partition passing the closure
test can have independent representatives only if
\[
                          p+|\pi|-1=r.
\]
The representative set itself is successful, so
Theorem~\ref{thm:fpt-parity-rank} forces it to be a basis. Therefore, for
one output assignment the number of partitions that need consideration is
\[
                 \left\{\begin{matrix}q+1\\r-p+1\end{matrix}\right\}.
\]
A partition in this Stirling layer still has to pass the nonlinear closure
test and the graphic independence test. The distinguished block is counted
in the partition but contributes no representative and no graphic rank.

\begin{corollary}[A bounded deficit gives single-exponential dependence]
\label{cor:bounded-deficit}
Assume full raw parity. For each output assignment not rejected by the
cardinality test, write $\delta_o=r-p_o$. If all these deficits satisfy
$0\le\delta_o\le\Delta$, then the triple-unit algorithm has running time
\[
                       3^k(\Delta+1)^{2k}\operatorname{poly}(N).
\]
For fixed $\Delta$, this is $2^{O(k)}\operatorname{poly}(N)$.
The hypothesis concerns every non-rejected output assignment, not merely
one successful assignment. If all assignments are rejected, take $\Delta=0$.
\end{corollary}
\begin{proof}
A partition of $q+1$ objects into $\delta+1$ blocks can be encoded by
labeling its distinguished block zero and labeling its other blocks
canonically with $1,\ldots,\delta$. Giving labels to the $q$ ordinary
objects gives the bound
\[
              \left\{\begin{matrix}q+1\\\delta+1\end{matrix}\right\}
                              \le(\delta+1)^q.
\]
There are at most $3^k$ output assignments and $q\le2k$ prerequisites.
The polynomial work on rejected assignments is absorbed by the same
bound. The proof also gives
$(\prod_i|U_i|)(\Delta+1)^Q\operatorname{poly}(N)$ for general supports.
\end{proof}

This corollary is already reflected in the implemented Stirling-layer
search; it requires no additional algorithm. In particular, for
polynomial-size inputs and uniformly bounded deficit, $k=O(\log n)$
gives polynomial local time, while $k=O((\log n)^2)$ gives
$n^{O(\log n)}$ local time. Proving such a uniform deficit bound is a
separate geometric problem.

The weighted dummy-completion construction in
Theorem~\ref{thm:fpt-minimum-weight-tree} applies without change to
\eqref{eq:output-groups}. Thus the stronger output-only bound also applies
to exact minimum-weight successful-tree optimization. The supplied code
implements the two unweighted strategies; this weighted extension is
proved but not implemented.

An exact negative certificate covers every output assignment and every
required distinguished partition. A failed partition has either a failed
restricted-closure test, which the checker recomputes, or a Rado rank
defect for \eqref{eq:output-groups}. Canonical restricted-growth encodings
identify partitions without ambiguity. A checker verifies uniqueness and
the exact Bell or Stirling count, independently reconstructs each branch,
and authenticates raw parity before permitting the rank restriction.
Reproducing an optimizer or trusting a producer's branch log is unnecessary.

The logical closure test provides a second protection against circular
justification. Chosen outputs do not become initial seeds merely because
they are listed in the assignment. They must be derived from $U\cup S_\pi$
by the actual restricted rules.
